
%%% -------------------- Titulní strana --------------------
\begin{titlepage}
\begin{center}
\large \vspace*{\fill}
\thispagestyle{empty}

\LARGE

{ \huge \textbf{Gymnázium Arabská, Praha 6, Arabská 14}}

{\LARGE Obor programování }

\vfill
\includegraphics{logogyarab.png}
\vspace{15pt}

\vfill

{\huge \textbf{Audiosociální síť: Ding}}

\vfill

Rory Beneš, Kryštof Breburda, Adam Suchý

\vfill

{\large Květen, 2023}

\vspace*{\fill}
\end{center}
\end{titlepage}

%%% -------------------- Prohlášení --------------------
\thispagestyle{empty}
\addtocounter{page}{-1}
\vspace*{\fill}
Prohlašujeme, že jsme jedinými autory tohoto projektu, všechny citace jsou řádně označené a všechna 
použitá literatura a další zdroje jsou v práci uvedené. Tímto dle zákona 121/2000 Sb. (tzv. Autorský zákon) 
ve znění pozdějších předpisů udělujeme bezúplatně škole Gymnázium, Praha 6, Arabská 14 oprávnění k výkonu 
práva na rozmnožování díla (§ 13) a práva na sdělování díla veřejnosti (§ 18) na dobu časově neomezenou a 
bez omezení územního rozsahu.

\vspace{2cm}
V .................. dne .......................

\rightline{Rory Beneš \hspace{0.5cm} ............................ \hspace{2cm}}

\vspace{2cm}

\rightline{Kryštof Breburda \hspace{0.5cm} ............................ \hspace{2cm}}

\vspace{2cm}

\rightline{Adam Suchý \hspace{0.5cm} ............................ \hspace{2cm}}

\vspace{2cm}

%%% -------------------- Anotace --------------------
\newpage
\begin{abstract}
	V rámci tohoto ročníkového projektu jsme vytvořili sociální síť Ding, která je vyhrazena pro zvukové
	nahrávky bez mluveného slova, filtrovaného umělou inteligencí. Ding je webová stránka napsaná v 
	programovacím jazyce JavaScript ve frameworku Vue.js. Serverová část sociální sítě je napsána v jazyce 
	Rust a využívá databázi PostgreSQL.
\end{abstract}

%%% -------------------- Obsah --------------------
\tableofcontents
